\section*{Chapitre \ref{ch:nw:the-vendor-map} --- La carte des fournisseurs}

\begin{solution}{exo:nw:the-vendor-map:1}
Il reçoit les flux une seule fois, les décode et les normalise, tient les carnets d'ordres et les diffuse à l'intérieur de la société. Centraliser évite un double décodage et des serveurs en double, au prix
d'un réseau de distribution et de sa latence.
\end{solution}

\begin{solution}{exo:nw:the-vendor-map:2}
Une société qui vend des logiciels de trading ou de post-marché aux sociétés de trading et les connecte aux bourses pour leur compte. La bourse la certifie
afin que son logiciel se comporte correctement sur les interfaces de la bourse, ce qui protège la bourse et tous les autres participants.
\end{solution}

\begin{solution}{exo:nw:the-vendor-map:3}
Elles ne se trouvent pas sur le chemin qui va des données aux ordres, donc leur panne n'arrête pas le trading. Mais chaque horodatage exigé par les régulateurs
dépend d'elles, si bien qu'une panne longue laisse la société incapable de prouver ses registres.
\end{solution}

\begin{solution}{exo:nw:the-vendor-map:4}
Le \omterm{def:nw:the-vendor-map:tickerplant}{ticker plant} dispose d'un chemin parallèle passant par les feed handlers internes, si bien que son retrait laisse subsister un chemin des données aux ordres.
\end{solution}

\begin{solution}{exo:nw:the-vendor-map:5}
1\,782~: les dépenses d'Equinix passent à 1.44~millions de dollars sur 3.91~millions de dollars, et son carré de part domine l'indice.
\end{solution}

\begin{solution}{exo:nw:the-vendor-map:6}
Un sous-fournisseur commun~: deux fournisseurs de serveurs qui utilisent les mêmes \omterm{def:nw:networking-for-trading:frame}{cartes réseau}, le même firmware ou les mêmes puces tombent en panne ensemble pour un même défaut. La carte
devrait consigner les fabricants des pièces, et non seulement les vendeurs.
\end{solution}

\begin{solution}{exo:nw:the-vendor-map:7}
AMD et Blackcore~; Equinix sort de la liste, de même que la \omterm{def:nw:colocation-products-and-how-they-are-sold:colo}{colocation} et les \omterm{def:nw:colocation-products-and-how-they-are-sold:mmr}{interconnexions} issues des composants uniques, et l'indice tombe à 1\,104.
\end{solution}

\begin{solution}{exo:nw:the-vendor-map:8}
La part de dépense n'est pas une dépendance~: un fournisseur à faible part peut fournir un composant situé sur le seul chemin qui va des données aux ordres, et le
plus gros fournisseur ici est aussi un point de défaillance dont la sortie prend un semestre.
\end{solution}

\begin{solution}{pb:nw:the-vendor-map:1}
\textbf{Partie I.}
\begin{enumerate}
\item La \omterm{def:nw:colocation-products-and-how-they-are-sold:colo}{colocation} et les \omterm{def:nw:colocation-products-and-how-they-are-sold:mmr}{interconnexions} (cross-connect) chez Equinix, une \omterm{def:nw:switches-and-layer-1-devices:l1}{duplication de couche 1} (fan-out) chez Arista, des commutateurs chez Cisco, un \omterm{def:nw:the-vendor-map:tickerplant}{ticker plant} chez Exegy, des serveurs
  (Blackcore), des \omterm{def:nw:networking-for-trading:frame}{cartes réseau} (AMD), la gestion des ordres (Trading Technologies), des \omterm{def:nw:time-synchronisation:gm}{horloges grand maître} (Meinberg), des outils d'analyse (Pico), un faisceau hertzien
  (Quincy Data), et ses propres feed handlers et sa passerelle.
\item La liste des fournisseurs certifiés d'une plateforme de négociation, les pages produits des fournisseurs et les lignes des chapitres précédents.
\item 3.19~millions de dollars~; Equinix 22.6~\%, Quincy Data 15.0~\%, travail interne 14.1~\%, Exegy 12.5~\%.
\item 1\,149.
\end{enumerate}
\textbf{Partie II.}
\begin{enumerate}[resume]
\item La \omterm{def:nw:colocation-products-and-how-they-are-sold:colo}{colocation}, les \omterm{def:nw:colocation-products-and-how-they-are-sold:mmr}{interconnexions}, les serveurs de stratégie, leurs \omterm{def:nw:networking-for-trading:frame}{cartes réseau} et la passerelle.
\item AMD, Blackcore et Equinix.
\item Chacun a un composant parallèle provenant d'un autre fournisseur ou développé en interne.
\item 26, 12, 6, 4 et 0 semaines~: 48 au total.
\end{enumerate}
\textbf{Partie III.}
\begin{enumerate}[resume]
\item L'indice tombe à 1\,119 et Blackcore sort des points fournisseurs~; les serveurs restent un composant unique.
\item Un deuxième fabricant de \omterm{def:nw:networking-for-trading:frame}{cartes réseau} sur certains serveurs, ou des cartes d'une autre marque dans les machines du second fournisseur de serveurs.
\item Le site de \omterm{def:nw:colocation-products-and-how-they-are-sold:colo}{colocation} est celui de la bourse~; un second site signifie un bâtiment d'un autre opérateur, ses propres \omterm{def:nw:colocation-products-and-how-they-are-sold:mmr}{interconnexions} et un
  chemin plus long, et un semestre pour déménager.
\item Rien pour le trading automatisé~; les flux manuels et d'agence s'arrêtent, et les registres et la supervision se dégradent.
\end{enumerate}
\textbf{Partie IV.}
\begin{enumerate}[resume]
\item \emph{Résultat nommé}~: la pile de la société exemple a un indice de concentration de 1\,149 sur les dépenses, et cinq composants dont le fournisseur unique
  (ou la société elle-même) interrompt le trading~: la \omterm{def:nw:colocation-products-and-how-they-are-sold:colo}{colocation} et les \omterm{def:nw:colocation-products-and-how-they-are-sold:mmr}{interconnexions} (Equinix), les serveurs de stratégie (Blackcore), leurs \omterm{def:nw:networking-for-trading:frame}{cartes réseau} (AMD)
  et la passerelle interne~; les points fournisseurs sont AMD, Blackcore et Equinix.
\item Publié~: les fournisseurs et ce qu'ils vendent. Pour la société~: sa pile, son graphe, ses dépenses et ses délais de sortie.
\item L'indice compte l'argent~; les points fournisseurs comptent les chemins. Un composant bon marché sur le seul chemin compte davantage qu'un composant coûteux
  disposant d'une pièce de rechange.
\item Pour chaque point fournisseur~: comment la société continue de trader ou s'arrête sans risque, les pièces de rechange, le second site, et qui appelle qui.
\item Ce qu'il faudrait pour partir, en combien de temps, et si quelque chose a changé qui fait du fournisseur un point de défaillance.
\item Le fournisseur du \omterm{def:nw:the-vendor-map:tickerplant}{ticker plant} disparaît et le travail en interne augmente~; la société assume désormais les pannes de ce composant, et sa sortie devient la sienne.
\item Dans la nomenclature (bill of materials) et le registre des risques du plan de connectivité.
\item Être le seul fournisseur d'un composant situé sur le seul chemin qui va des données aux ordres.
\end{enumerate}
\end{solution}

\begin{solution}{iq:nw:the-vendor-map:1}
Un \omterm{def:nw:the-vendor-map:tickerplant}{ticker plant} décode une fois et distribue, ce qui économise des serveurs et un travail en double, au prix d'un saut réseau et d'un point central
de défaillance~; les feed handlers embarqués décodent dans chaque stratégie, plus vite et de manière indépendante, au prix d'un travail dupliqué et de nombreuses copies à tenir
à jour.
\iqlookfor{Duplication face à la latence~; distribution~; domaines de défaillance.}
\end{solution}

\begin{solution}{iq:nw:the-vendor-map:2}
On dessine le chemin qui mène des données de marché aux ordres sous la forme d'un graphe comportant chaque composant et son fournisseur, on retire tour à tour chaque composant et chaque fournisseur,
et l'on vérifie si un chemin subsiste~; on teste ensuite les réponses lors d'exercices de bascule (failover).
\iqlookfor{Graphe~; coupures par composant et par fournisseur~; tests.}
\end{solution}

\begin{solution}{iq:nw:the-vendor-map:3}
On chiffre les options~: une mise à niveau avec le fournisseur, un changement de fournisseur ou une construction interne~; on prévoit la certification et les tests que chacune exige, et
on lance dès maintenant celle dont le délai est le plus long.
\iqlookfor{Délai de sortie~; certification~; délai d'approvisionnement.}
\end{solution}

\begin{solution}{iq:nw:the-vendor-map:4}
Pour les pièces de rechange, les pilotes, les compétences et une performance constante~; pas lorsque ce fournisseur devient le seul fournisseur des serveurs du chemin de trading,
où un second fournisseur pour certaines machines supprime un point de défaillance.
\iqlookfor{Standardisation face à la dépendance.}
\end{solution}

\begin{solution}{iq:nw:the-vendor-map:5}
Les stratégies automatisées, dotées de leurs propres passerelles, continuent de trader~; les ordres manuels et les allocations s'arrêtent. Il faut passer par le courtier ou par les outils propres
à la bourse pour les ordres manuels urgents, puis réconcilier les positions à partir des drop copies.
\iqlookfor{Ce qui se trouve sur le chemin~; solutions de repli~; réconciliation.}
\end{solution}

\begin{solution}{iq:nw:the-vendor-map:6}
On achète les couches banalisées (\omterm{def:nw:colocation-products-and-how-they-are-sold:colo}{colocation}, matériel, logiciels de connectivité certifiés), on construit ce qui différencie (stratégies, passerelle
d'exécution), on double les sources des composants du chemin de trading lorsque c'est peu coûteux, on enregistre chaque fournisseur avec son délai de sortie, et on revoit la carte
à chaque renouvellement.
\iqlookfor{Achat contre construction~; chemin de trading~; sorties.}
\end{solution}
